\documentclass[14pt,dvipdfmx]{jsarticle}
\input{lecture-preamble.tex}

\begin{document}

\lectureheader{数学Ⅱ}{6}{微分法と積分法}{1}{微分係数と導関数}{2-A}{導関数}
  {182,183}{導関数を理解し, 定義と公式を使って求めよう}

% 1-1-CD の定義 6.1.3、6.1.4 の続き
\setcounter{definition}{4}

% ============ 1ページ目（表・左） ============
\vspace{-2.0zw}%
\heading{〜前時の復習〜}
\vspace{-0.35em}
\begin{reviewbox}{微分係数と接線}
  \begin{enumerate}[label=\textbf{\arabic*}, leftmargin=1.5zw, itemsep=0.2em]
    \item $x=a$ における微分係数は
          \par\vspace{0.3em}
          {\renewcommand{\fitblankpad}{3zw}%
          \fitblankbf{f'(a)=\lim_{h \to 0}\dfrac{f(a+h)-f(a)}{h}}}
          \mbox{である.}
    \item 接線の傾きは, 微分係数 \fitblankbf{f'(a)} に等しい.
  \end{enumerate}
\end{reviewbox}

\vspace{0.45em}
\heading{微分係数を $a$ の式で表す}
\vspace{-0.25em}
関数 $f(x)$ が与えられたとき, $x=a$ における微分係数 $f'(a)$ を計算すると,
$f'(a)$ が $a$ の式で表される.
関数 $f(x)=x^{2}$ で確かめてみよう.
\par\vspace{0.15em}
\[
\begin{aligned}
  f'(a)
  &=\lim_{h \to 0}\frac{f(a+h)-f(a)}{h}
   =\lim_{h \to 0}\frac{(a+h)^{2}-a^{2}}{h} \\[0.4em]
  &=\lim_{h \to 0}\frac{2ah+h^{2}}{h}
   =\lim_{h \to 0}\frac{h(2a+h)}{h}
   =\lim_{h \to 0}(2a+h)=\fitblankbf{2a}
\end{aligned}
\]
\vspace{-0.4em}
すなわち $f'(a)=\fitblankbf{2a}$ ……①
\par\vspace{0.45em}
①を利用すると, 極限を計算し直さなくても, $x$ のいろいろな値における微分係数が求められる.
たとえば $f'(3)$ は, ①に $a=3$ を代入して
$f'(3)=\fitblankbf{2\cdot 3=6}$ である.

\vfill
\nextpage

% ============ 2ページ目（表・右） ============
\vspace{-1.4zw}%
\begin{practiceprob}{5}
  関数 $f(x)=x^{2}$ について, ①を利用して, 次の微分係数を求めよ。
  \begin{tasks}[label=(\arabic*), label-width=1.6zw, item-indent=2zw, after-item-skip=0.12em](3)
    \task $f'(4)$
    \task $f'(0)$
    \task $f'(-2)$
  \end{tasks}
\end{practiceprob}
\vspace{-0.25em}
\begin{answerbox}{(1)}[7]
  ①より $f'(4)=2\cdot 4=8$
  \boxanswer{$8$}
\end{answerbox}
\begin{answerbox}{(2)}[7]
  ①より $f'(0)=2\cdot 0=0$
  \boxanswer{$0$}
\end{answerbox}
\begin{answerbox}{(3)}[7]
  ①より $f'(-2)=2\cdot(-2)=-4$
  \boxanswer{$-4$}
\end{answerbox}

\vspace{0.15em}
①で $a$ の値を $1$ つ決めると, $f'(a)$ の値もただ $1$ つ決まる.
すなわち $a$ を変数とみると, $f'(a)$ は $a$ の関数である.
\par\vspace{0.25em}
一般に, 関数 $f(x)$ において, $x$ の各値に微分係数を対応させると,
$x$ の新しい関数が得られる.
たとえば $f(x)=x^{2}$ では, $f'(x)=\fitblankbf{2x}$ である.

\vspace{0.25em}
\begin{definitionbox}{導関数}
  このようにして得られる関数を, もとの関数 $f(x)$ の\fitblankbf{\text{導関数}}といい,
  \fitblankbf{f'(x)} で表す.
  \[
    f'(x)=\fitblankbf{\lim_{h \to 0}\dfrac{f(x+h)-f(x)}{h}}
  \]
\end{definitionbox}

\vfill
\nextpage

% ============ 3ページ目（裏・左） ============
\vspace{-1.4zw}%
\noindent{\bfseries 例\,5}\quad
導関数の定義にしたがって, 次の関数の導関数を求める。

\vspace{0.1em}
\noindent (1)\ 関数 $f(x)=x$
\vspace{-0.15em}
\begin{answerbox}{}[7]
  $f'(x)=\lim_{h \to 0}\dfrac{(x+h)-x}{h}
    =\lim_{h \to 0}\dfrac{h}{h}=\lim_{h \to 0}1=1$
  \boxanswer{$1$}
\end{answerbox}

\noindent (2)\ 関数 $f(x)=x^{3}$
\vspace{-0.15em}
\begin{answerbox}{}[13]
  $(x+h)^{3}=x^{3}+3x^{2}h+3xh^{2}+h^{3}$ より\\[0.35em]
  $f'(x)=\lim_{h \to 0}\dfrac{(x+h)^{3}-x^{3}}{h}
    =\lim_{h \to 0}\dfrac{3x^{2}h+3xh^{2}+h^{3}}{h}$\\[0.45em]
  $=\lim_{h \to 0}(3x^{2}+3xh+h^{2})=3x^{2}$\\[0.15em]
  {\footnotesize $3xh$ と $h^{2}$ は, どちらも $0$ に限りなく近づく.}
  \boxanswer{$3x^{2}$}
\end{answerbox}

\noindent (3)\ 関数 $f(x)=2$
\vspace{-0.15em}
\begin{answerbox}{}[6]
  $f'(x)=\lim_{h \to 0}\dfrac{2-2}{h}=\lim_{h \to 0}0=0$
  \boxanswer{$0$}
\end{answerbox}

\vspace{0.1em}
\textbf{注意}\quad
$c$ を定数とするとき, $f(x)=c$ の形の関数を\fitblankbf{\text{定数関数}}という.
定数関数の導関数は \fitblankbf{0} である.

\vspace{0.35em}
導関数は $y'$, $\dfrac{dy}{dx}$, $\dfrac{d}{dx}f(x)$ や $(x^{3})'$ とも書く.
これまでに $(x)'=1$, $(x^{2})'=2x$, $(x^{3})'=3x^{2}$, $(2)'=0$ である.
一般に, $n$ を正の整数とすると, 次の公式が成り立つ.

\vspace{-0.15em}
\begin{formulabox}{$x^{n}$ と定数関数の導関数}
  $(x^{n})'=\fitblankbf{nx^{n-1}}$,\qquad
  $(c)'=\fitblankbf{0}$
  \par\vspace{0.1em}
  {\footnotesize $c$ は定数とする. 証明は教科書 p.\,187.}
\end{formulabox}

\vfill
\nextpage

% ============ 4ページ目（裏・右） ============
\vspace{-1.4zw}%
\begin{practiceprob}{6}
  導関数の定義にしたがって, 次の関数の導関数を求めよ。
  \begin{tasks}[label=(\arabic*), label-width=1.6zw, item-indent=2zw, after-item-skip=0.08em](3)
    \task $f(x)=3x$
    \task $f(x)=-x^{2}$
    \task $f(x)=4$
  \end{tasks}
\end{practiceprob}
\vspace{-0.35em}
\begin{answerbox}{(1)}[6]
  $f'(x)=\lim_{h \to 0}\dfrac{3(x+h)-3x}{h}
    =\lim_{h \to 0}\dfrac{3h}{h}=3$
  \boxanswer{$3$}
\end{answerbox}
\begin{answerbox}{(2)}[9]
  $f'(x)=\lim_{h \to 0}\dfrac{-(x+h)^{2}-(-x^{2})}{h}$\\[0.35em]
  $=\lim_{h \to 0}\dfrac{-2xh-h^{2}}{h}
    =\lim_{h \to 0}(-2x-h)=-2x$
  \boxanswer{$-2x$}
\end{answerbox}
\begin{answerbox}{(3)}[6]
  $f'(x)=\lim_{h \to 0}\dfrac{4-4}{h}=\lim_{h \to 0}0=0$
  \boxanswer{$0$}
\end{answerbox}

\vspace{-0.15em}
\begin{practiceprob}{7}
  公式を用いて, 次の関数の導関数を求めよ。
  \begin{tasks}[label=(\arabic*), label-width=1.6zw, item-indent=2zw, after-item-skip=0.05em](2)
    \task $y=x^{4}$
    \task $y=x^{5}$
  \end{tasks}
\end{practiceprob}
\vspace{-0.45em}
\begin{answerbox}{(1)}[5]
  $(x^{4})'=4x^{3}$
  \boxanswer{$4x^{3}$}
\end{answerbox}
\begin{answerbox}{(2)}[5]
  $(x^{5})'=5x^{4}$
  \boxanswer{$5x^{4}$}
\end{answerbox}

\vspace{0.05em}
\begin{summarybox}
  \begin{enumerate}[label=\textbf{\arabic*}, leftmargin=1.5zw, itemsep=0.25em]
    \item 導関数は
          {\renewcommand{\fitblankpad}{2.2zw}%
          \fitblankbf{f'(x)=\lim_{h \to 0}\dfrac{f(x+h)-f(x)}{h}}}
          \mbox{である.}
    \item $(x^{n})'=\fitblankbf{nx^{n-1}}$,\ 
          $(c)'=\fitblankbf{0}$\mbox{である.}
  \end{enumerate}
\end{summarybox}

\vspace{0.05em}
\nexttime{和や定数倍の導関数を使って, 整式を微分してみよう!}

\end{document}
